% =====================================================================
%  Beneath the Mesa -- ENG 220 Project Report (Team 23)
%  Two-column report template styled to match the project proposal:
%  Palatino (URW Palladio) text, navy headings, blue rules and header.
%
%  Build:  pdflatex report -> bibtex report -> pdflatex report (x2)
%          (or `latexmk -pdf report.tex`; Overleaf does this automatically)
%
%  Orange [notes] mark what still needs to be written.
%  Dashed boxes mark figures that still need to be made.
% =====================================================================
\documentclass[10pt,twocolumn,letterpaper]{article}

% ---------------------------------------------------------------------
%  Switches
% ---------------------------------------------------------------------
\newif\ifdraftnotes \draftnotestrue   % \draftnotesfalse hides every [orange note]
\newif\ifcoverpage  \coverpagefalse   % \coverpagetrue adds a proposal-style cover page
\newif\ifroughdraft \roughdrafttrue    % \roughdraftfalse removes the red ROUGH DRAFT stamp

% ---------------------------------------------------------------------
%  Fonts and text
% ---------------------------------------------------------------------
\usepackage[T1]{fontenc}
\usepackage[utf8]{inputenc}
\usepackage{mathpazo}                 % URW Palladio, the proposal's typeface
\usepackage[protrusion=true,expansion=true]{microtype}
\usepackage{amsmath}

% ---------------------------------------------------------------------
%  Page layout
% ---------------------------------------------------------------------
\usepackage[letterpaper,top=0.8in,bottom=0.85in,left=0.75in,right=0.75in,
            columnsep=0.28in,headheight=20pt,headsep=8pt,footskip=26pt]{geometry}
\setlength{\parindent}{1.2em}
\setlength{\parskip}{0pt}

% ---------------------------------------------------------------------
%  Colors (sampled from the proposal PDF)
% ---------------------------------------------------------------------
\usepackage{xcolor}
\definecolor{titleblue}{RGB}{23,54,93}     % #17365D  headings, title
\definecolor{accentblue}{RGB}{54,95,145}   % #365F91  rules, header, links
\definecolor{tintblue}{RGB}{237,241,247}   % light panel fill
\definecolor{draftorange}{RGB}{190,84,20}  % draft notes only
\definecolor{draftred}{RGB}{192,0,0}       % ROUGH DRAFT stamp

% ---------------------------------------------------------------------
%  Headings
% ---------------------------------------------------------------------
\usepackage{titlesec}
\usepackage{needspace}
\usepackage{etoolbox}
\titleformat{\section}{\color{titleblue}\large\bfseries}{\thesection}{0.55em}{}
  [{\color{accentblue}\titlerule[0.6pt]}]
\titlespacing*{\section}{0pt}{2.0ex plus 0.6ex minus 0.3ex}{1.0ex}
\titleformat{\subsection}{\color{accentblue}\normalsize\bfseries}{\thesubsection}{0.55em}{}
\titlespacing*{\subsection}{0pt}{1.4ex plus 0.4ex minus 0.2ex}{0.5ex}
\titleformat{\subsubsection}[runin]{\normalsize\bfseries\itshape}{}{0pt}{}[.]
\titlespacing*{\subsubsection}{0pt}{1.0ex plus 0.3ex}{0.6em}
\pretocmd{\section}{\Needspace{4\baselineskip}}{}{}       % keep headings with their text
\pretocmd{\subsection}{\Needspace{3\baselineskip}}{}{}

% ---------------------------------------------------------------------
%  Running header and footer (same as the proposal)
% ---------------------------------------------------------------------
\usepackage{fancyhdr}
\pagestyle{fancy}
\fancyhf{}
\fancyhead[L]{\small\color{accentblue}Project Report}
\fancyhead[R]{\small\color{accentblue}Aquifer \& Well Resources: Sandoval County, NM}
\fancyfoot[C]{\small\thepage}
\renewcommand{\headrulewidth}{0.4pt}
\renewcommand{\headrule}{\hbox to\headwidth{\color{accentblue}\leaders\hrule height\headrulewidth\hfill}}

\fancypagestyle{firstpage}{%
  \fancyhf{}%
  \renewcommand{\headrulewidth}{0pt}%
  \fancyfoot[L]{\footnotesize\color{accentblue}ENG 220 Project Report \,\textbullet\, Team 23 \,\textbullet\, \reportstatus}%
  \fancyfoot[R]{\small\thepage}%
}

% ---------------------------------------------------------------------
%  Lists, tables, figures, captions
% ---------------------------------------------------------------------
\usepackage{enumitem}
\setlist{itemsep=1pt,topsep=3pt,parsep=0pt}
\setlist[itemize]{leftmargin=1.2em,label={\color{accentblue}\textbullet}}
\usepackage{booktabs,tabularx,array}
\renewcommand{\arraystretch}{1.12}
\usepackage{graphicx}
\graphicspath{{figures/}}
\usepackage{caption}
\captionsetup{font=small,labelfont={bf,color=titleblue},labelsep=period,justification=justified}
\captionsetup[table]{position=top,skip=4pt}  % for tables you add later
\usepackage{tikz}
\usepackage{stfloats}                 % lets full-width figures sit at the page bottom too

% ---------------------------------------------------------------------
%  References and links
% ---------------------------------------------------------------------
\usepackage[round,authoryear]{natbib}
\setlength{\bibsep}{2.5pt}
\renewcommand{\bibfont}{\small\raggedright}
\usepackage{xurl}
\usepackage[colorlinks=true,linkcolor=accentblue,citecolor=accentblue,urlcolor=accentblue,
            pdftitle={Beneath the Mesa: Aquifer and Well Water Resources Across Sandoval County, New Mexico},
            pdfauthor={Djordjevic, Gonzales, Wauneka, Watkins, Yousef}]{hyperref}
\urlstyle{same}

% ---------------------------------------------------------------------
%  Draft helpers
% ---------------------------------------------------------------------
\newcommand{\reportstatus}{\ifroughdraft Rough draft, November 2, 2026\else Final, November 16, 2026\fi}
% Red stamp shown on page 1 and the cover page while \roughdrafttrue
\newcommand{\draftstamp}[1][\large]{\ifroughdraft%
  \tikz\node[draw=draftred,line width=1.1pt,rounded corners=2pt,inner xsep=9pt,inner ysep=4pt,
    text=draftred,font=#1\bfseries]{\textls[180]{ROUGH DRAFT}};\fi}
\newcommand{\draftnote}[1]{\ifdraftnotes{\color{draftorange}\itshape[#1]}\fi}
% Dashed box that stands in for a figure that does not exist yet.
%   \figplaceholder[height]{what the figure will show}
% When the figure is ready, save it in figures/ and replace the
% placeholder with:  \includegraphics[width=\linewidth]{file-name}
\newcommand{\figplaceholder}[2][1.8in]{%
  \begin{tikzpicture}
    \draw[accentblue!70,dashed,line width=0.6pt,fill=tintblue,rounded corners=2pt]
      (0,0) rectangle (\linewidth,#1);
    \node[text width=0.82\linewidth,align=center,font=\small\itshape,text=accentblue]
      at (0.5\linewidth,0.5*#1) {#2};
  \end{tikzpicture}}
\newcommand{\infohead}[1]{{\color{titleblue}\footnotesize\bfseries\MakeUppercase{#1}}\par\vspace{2pt}}

% =====================================================================
\begin{document}

% ---------------------------------------------------------------------
%  Optional cover page (same layout as the proposal cover)
% ---------------------------------------------------------------------
\ifcoverpage
\begin{titlepage}
  \centering
  \vspace*{1.6in}
  {\color{accentblue}\rule{\textwidth}{0.8pt}}\par\vspace{1.4em}
  {\color{titleblue}\fontsize{24}{29}\selectfont\bfseries Beneath the Mesa\par}
  \vspace{0.9em}
  {\color{accentblue}\fontsize{17}{22}\selectfont Assessing Aquifer and Well Water Resources\\
    Across Sandoval County, New Mexico\par}
  \vspace{1.3em}
  {\color{accentblue}\rule{\textwidth}{0.8pt}}\par\vspace{1.6em}
  {\itshape A County-Scale Screening Study with Cuba, New Mexico as a Focus Area\par}
  \vspace{1.3in}
  {\color{titleblue}\large\bfseries Project Report\par}
  \ifroughdraft\vspace{1.2em}\draftstamp[\Large]\par\vspace{-0.4in}\fi
  \vspace{1.2in}
  \begin{tabular}{@{}r@{\hspace{0.8em}}l@{}}
    \textbf{Course:}     & ENG 220 \\
    \textbf{Instructor:} & Ramiro Jordan \\
    \textbf{Team:}       & \#23 Petar Djordjevic, Dylan Gonzales, James Wauneka,\\
                         & \phantom{\#23 }Micah Watkins, Meckrem Yousef \\
    \textbf{Date:}       & \ifroughdraft November 2, 2026 (rough draft)\else November 16, 2026\fi \\
  \end{tabular}
  \vfill
\end{titlepage}
\fi

% ---------------------------------------------------------------------
%  Title block and abstract (full width)
% ---------------------------------------------------------------------
\twocolumn[{%
\begin{center}
  \ifroughdraft\vspace*{-1.5ex}\draftstamp\par\vspace{1.4ex}\fi
  {\color{accentblue}\rule{\textwidth}{0.8pt}}\par\vspace{0.9ex}
  {\color{titleblue}\fontsize{22}{26}\selectfont\bfseries Beneath the Mesa\par}
  \vspace{0.7ex}
  {\color{accentblue}\fontsize{14.5}{18}\selectfont Assessing Aquifer and Well Water Resources\\
    Across Sandoval County, New Mexico\par}
  \vspace{0.9ex}
  {\color{accentblue}\rule{\textwidth}{0.8pt}}\par\vspace{0.9ex}
  {\itshape A County-Scale Screening Study with Cuba, New Mexico as a Focus Area\par}
  \vspace{2.0ex}
  {\large Petar Djordjevic, Dylan Gonzales, James Wauneka, Micah Watkins, and Meckrem Yousef\par}
  \vspace{0.6ex}
  {\small ENG 220, Team 23 \,\textbullet\, Instructor: Ramiro Jordan\par}
\end{center}
\vspace{1.2ex}
{\color{accentblue}\rule{\textwidth}{0.4pt}}\par\vspace{1.4ex}
\noindent
\begin{minipage}[t]{0.25\textwidth}
  \raggedright\footnotesize
  \infohead{Report info}
  \ifroughdraft Status: {\color{draftred}\bfseries Rough draft}\\
  Date: November 2, 2026\\
  \else Status: Final\\
  \fi Final due: November 16, 2026\par\vspace{1.4ex}
  \infohead{Keywords}
  Groundwater\\ Aquifers\\ Wells\\ Sandoval County\\ Cuba, New Mexico\par\vspace{1.4ex}
  \infohead{Online}
  \href{https://eng220aquiferproject.streamlit.app/}{Interactive dashboard}\\
  \href{https://github.com/pdjordj11/eng_220_aquifer_project}{Code and data (GitHub)}\\
  \href{https://pdjordj11.github.io/eng_220_aquifer_project/}{Project documentation}
\end{minipage}%
\hfill{\color{accentblue}\vrule width 0.4pt}\hfill
\begin{minipage}[t]{0.70\textwidth}
  \infohead{Abstract}
  \small
  Water scarcity is a persistent challenge across New Mexico. Sandoval County spans several
  hydrogeologic settings, and its communities depend on a mix of groundwater wells, surface water,
  and acequia irrigation systems. This project focuses on groundwater: it evaluates the condition
  of the county's aquifers and wells, with Cuba, New Mexico, on the eastern edge of the San Juan
  Basin, as a focused case study. The project is motivated by concern about declining groundwater
  near Cuba, but rather than assume how severe that decline is or what is causing it, we test what
  the available public data actually show, and where they are too sparse to say.
  \draftnote{Add the data used, the main findings, and what they mean, in three to four
  sentences, once the analysis is complete.}
\end{minipage}
\par\vspace{1.6ex}
{\color{accentblue}\rule{\textwidth}{0.4pt}}\par\vspace{2.6ex}
}]
\thispagestyle{firstpage}

% =====================================================================
\section{Introduction}
% =====================================================================
\subsection{Background and Motivation}
\draftnote{Background on water resources in New Mexico and Sandoval County, and the concern about
declining groundwater near Cuba that motivated the project.}

\subsection{Approach}
The study works at two levels:
\begin{itemize}
  \item \textbf{County-wide screening.} Map and compare aquifers, wells, and groundwater levels
        across Sandoval County. Because the county crosses several geologic regions, comparisons
        are grouped by aquifer unit instead of treating the county as one uniform aquifer.
  \item \textbf{Cuba case study.} A closer look at the Ojo Alamo Sandstone and San Jose Formation
        near Cuba, using published maps of both formations, aquifer-test records from nearby
        wells, and stream measurements on the Rio Puerco at Cuba.
\end{itemize}

\subsection{Research Questions}
\begin{enumerate}[label={\color{titleblue}\bfseries RQ\arabic*.},leftmargin=2.6em,labelsep=0.4em]
  \item How have groundwater levels changed at monitored wells in Sandoval County, and which wells
        show sustained decline, stability, or recovery?
  \item What do well records show about well depth, depth to water, and the aquifer units that
        wells in different parts of the county draw from?
  \item What do available records show about groundwater near Cuba, and are they enough to
        evaluate the concern that its main aquifer is approaching severe depletion?
  \item How do the Ojo Alamo Sandstone and San Jose Formation compare near Cuba in extent,
        thickness, and hydraulic properties?
  \item Where is monitoring coverage strong, and where are the most important data gaps?
  \item Which indicators can be combined into a dashboard that shows groundwater stress without
        overstating how certain the data are?
\end{enumerate}
The project proposal also asked about groundwater quality, acequias, and links between
groundwater and surface water. Those questions were dropped because the data needed to answer
them were not available (Section~\ref{sec:limits}).

% ---------------------------------------------------------------------
%  Figure 1: county map (LaTeX places it at the top of the next page)
% ---------------------------------------------------------------------
\begin{figure*}[t]
  \centering
  \figplaceholder[2.2in]{Figure 1 (to be added by the analysis lead): county map showing
  water-level monitoring wells, aquifer-test wells colored by formation, and the Rio Puerco gage,
  with Cuba marked.}
  \caption{\draftnote{Caption to be written.}}
  \label{fig:map}
\end{figure*}

% =====================================================================
\section{Study Area}\label{sec:study}
% =====================================================================
\subsection{Sandoval County}
\draftnote{Describe the county, its hydrogeologic settings, and its main communities and water
sources (Fig.~\ref{fig:map}).}

\subsection{Cuba Focus Area}
Cuba lies on the eastern edge of the San Juan Basin. The case study focuses on the Ojo Alamo
Sandstone and San Jose Formation near Cuba \citep{usgs_rasa_sjb,kernodle1996}.
\draftnote{Describe Cuba, its water supply, and the aquifers it draws from, and define the
boundaries of the focus area.}

% =====================================================================
\section{Data}\label{sec:data}
% =====================================================================
Town-specific data for Cuba are limited, so the study relies on public datasets that can be
filtered to Sandoval County and the area around Cuba. Eight datasets were collected
(Table~\ref{tab:datasets}). For seven of them, a one-page brief records what the dataset contains,
which research questions it can help answer, which plots to make from it, and what it can't
show; a plotting guide compares all seven. The raw files are kept unedited in the project
repository \citep{team23_repo}.

% ---------------------------------------------------------------------
%  Table 1: datasets (full width; LaTeX places it at the top of a page)
% ---------------------------------------------------------------------
\begin{table*}[t]
  \caption{Datasets collected for the study. Distances are measured from Cuba.}
  \label{tab:datasets}
  \footnotesize
  \begin{tabularx}{\textwidth}{@{}>{\raggedright}p{0.9in}>{\raggedright}p{0.95in}
      >{\raggedright\arraybackslash}X>{\raggedright\arraybackslash}X l@{}}
    \toprule
    \textbf{Dataset} & \textbf{Source} & \textbf{Contents} & \textbf{Relevance to Cuba} & \textbf{Used for} \\
    \midrule
    Groundwater-level statistics & USGS & Monthly and yearly depth to water at 9 wells on 4 sites,
      1998--2026 & Wells are 40--59 miles away, in the eastern part of the county & RQ1, RQ5 \\
    Groundwater levels & NM Bureau of Geology & 456 depth-to-water measurements at 87 wells,
      1997--2019, with well depth & Nearest well is about 24 miles away & RQ1, RQ2, RQ5 \\
    Aquifer tests & NM Office of the State Engineer & 179 tests in Sandoval County, 1930--2008:
      well construction, formation, and hydraulic properties & Includes tests at a well about
      half a mile from Cuba & RQ2--RQ4 \\
    Wells with well logs & NM Office of the State Engineer & Drilling records for 6 Sandoval
      County wells, 1978--2022: borehole depth and depth to first water & Link to Cuba's wells
      not verified & RQ2, RQ5 \\
    Ojo Alamo Sandstone & USGS map atlas, 1990 & Maps of depth to the top of the sandstone, its
      thickness, and wells completed in it & Regional; covers the Cuba area & RQ4 \\
    San Jose Formation near Lindrith & NM Geological Society, 2025 & Water-bearing units,
      thickness of the Cuba Mesa Member, and groundwater levels & Centered near Lindrith, about
      20 miles north & RQ4 \\
    Rio Puerco at Cuba & USGS station 08332525 & 10 streamflow measurements (2019--2026) and
      4 gage-height readings & At Cuba, but measures the river, not groundwater & RQ3, RQ5 \\
    Aquifer sensitivity & NM Water Data & Sensitivity class for 613 mapped areas & Can't be
      located: the file has no coordinates & Not used \\
    \bottomrule
  \end{tabularx}
\end{table*}

\subsection{Datasets Used}
\subsubsection{Groundwater levels}
Two datasets measure depth to water below land surface, where a larger number means a lower
water table. The USGS statistics \citep{usgs_stats} give long monthly records for 9 wells at four
sites: Sandia, Lincoln Middle School in Rio Rancho, Bernalillo, and Dome Road. The same download also covers stream gages, canals, and
reservoirs, but only the groundwater records are used. The NM Bureau of Geology measurements
\citep{nmwd_levels} cover more wells but fewer years: most were measured in 1997--1998, and only
5 wells were measured again in 2016--2019, so few of them can show a long-term trend.

\subsubsection{Well and aquifer records}
The State Engineer's aquifer-test records \citep{ose_aquifertests} are the only dataset with
hydraulic properties, such as transmissivity and storativity, and they name the formation for
125 of the 179 tests. They include 8 Ojo Alamo Sandstone records at 5 locations, among them tests
from 1957 and 1960 at a school well about half a mile from Cuba, but only one San Jose Formation
well. The State Engineer's well logs \citep{ose_welllogs} add borehole depth and the depth at which
water was first found on the day each well was drilled, which is not a current water level.
\draftnote{Add the exact download page for the aquifer-test records.}

\subsubsection{Geology of the Cuba-area aquifers}
Two publications describe the formations studied for RQ4: a USGS map atlas of the Ojo Alamo
Sandstone across the San Juan Basin \citep{thorn1990}, and a study of the San Jose Formation near
Lindrith that identifies four water-bearing units and maps the Cuba Mesa Member
\citep{blandford2025}. Both are maps and reports rather than data tables, and neither shows which
formation Cuba's municipal wells draw from.

\subsubsection{Surface water at Cuba}
USGS station 08332525 on the Rio Puerco \citep{usgs_riopuerco} is the only dataset measured at
Cuba itself. Six of its 10 streamflow visits found no flow, and the other four measured 6.9 to
20.4 cubic feet per second. Because these are occasional visits rather than a continuous record,
the data are used only as context for the Cuba case study.

\subsubsection{Not used}
The aquifer sensitivity file \citep{nmwd_sensitivity} has no coordinates or shapes, so it can't be
mapped or matched to Cuba without its separate map layer, which was not obtained.

\subsection{Data Not Available}
The most important missing data are Cuba's municipal well records: well IDs, the formation each
well draws from, pumping volumes, and water-level history, held by the Village of Cuba or the
State Engineer \citep{nmose}. Every dataset brief identified these records as the missing piece for
RQ3. Also not obtained were the USGS digital aquifer data for the San Juan Basin
\citep{usgs_rasa_sjb}, groundwater-quality records such as the Brackish Water Assessment
\citep{nmbg_bwa}, acequia locations \citep{isc_acequias}, and records from other monitoring
networks \citep{ngwmn,cgmn}. Section~\ref{sec:limits} describes what their absence means for the
findings.

% =====================================================================
\section{Methods}\label{sec:methods}
% =====================================================================
\subsection{Data Cleaning}
Cleaning makes the data correct and consistent without drawing conclusions; averaging, trends,
and narrowing to the Cuba area happen in the analysis. Raw downloads are kept unedited. A MATLAB
script, \texttt{cleaning/clean\_csvs.m}, reads them and writes the cleaned files
(Table~\ref{tab:cleaned}), so anyone can rerun it and get the same result. As it runs, it prints
checks such as row counts, duplicate rows, and the list of formation names.

\begin{table}[t]
  \caption{Cleaned files produced so far.}
  \label{tab:cleaned}
  \footnotesize
  \begin{tabularx}{\linewidth}{@{}>{\raggedright}p{1.15in}r>{\raggedright\arraybackslash}X@{}}
    \toprule
    \textbf{File} & \textbf{Rows} & \textbf{Contents} \\
    \midrule
    \texttt{groundwater\_levels\_monthly.csv} & 2,673 & Monthly mean depth to water at
      9 wells, June 1998 to September 2026 \\
    \texttt{aquifer\_tests.csv} & 179 & One row per aquifer test in Sandoval County, 1930 to
      2008, with 17 columns \\
    \bottomrule
  \end{tabularx}
\end{table}

\subsubsection{Shared conventions}
Every cleaned file starts with the same six columns: site ID, site name, latitude, longitude,
date, and source. Each has one row per measurement, sorted by site and date. Numbers are stored
without units or commas, dates are written as YYYY-MM-DD, and column names end with their unit,
such as \texttt{well\_depth\_ft}. Missing values are left blank rather than filled with
placeholders such as 0 or $-9999$, which a chart would plot as real values.

\subsubsection{Groundwater levels}
From the USGS statistics, the script keeps the depth-to-groundwater rows that are monthly means
and renames the columns to the shared names. Wells with an A or B suffix are kept as separate
sites, because they share a location but are measured at different depths. The cleaned file was
checked against the raw data: it has the same rows, dates, and values, with no duplicate
site-and-date rows.

\subsubsection{Aquifer tests}
The script keeps all 179 Sandoval County tests, including repeat tests on the same well, and the
16 raw columns the analysis needs: well construction, formation, water level, pumping rate, and
hydraulic properties. Empty and administrative columns are dropped. It converts numbers stored as
text, standardizes dates, and cleans formation names by trimming spaces, converting them to
uppercase, and fixing a misspelling. Three decisions were needed:
\begin{itemize}
  \item One hydraulic conductivity was recorded as a range, 11 to 21 ft/day; the midpoint,
        16 ft/day, is used.
  \item 104 tests list January 1, 1900, as the date of the starting water level, meaning the
        date is unknown, so that column is left out.
  \item 30 tests are dated January 1, which may mean only the year is known; they are kept as is.
\end{itemize}
Many tests are missing values, such as the formation (54 tests) and storativity (165 tests); these
cells are left blank.

\subsubsection{Open issues}
Most of the State Engineer's well IDs in the aquifer-test file are incomplete: 157 of the 179 are
partial, so the cleaned file has 73 site IDs for about 149 distinct wells. The analysis therefore
tells wells apart by name and location rather than by ID. The Bureau of Geology water levels, the
Rio Puerco measurements, and the well logs still need cleaning; the first and last of these also
need their UTM coordinates converted to latitude and longitude.
\draftnote{Update this paragraph once the remaining files are cleaned.}

\subsection{Analysis}
The analysis focuses on measurable indicators, such as depth to water, water-level change over
time, well depth, aquifer unit, and hydraulic properties from aquifer tests, rather than a full
hydrologic model. Where records
are sparse, the results say so instead of filling the gaps.
\draftnote{Analysis lead: describe the method used for each research question.}

\subsection{Classifying Findings}
Each finding is labeled as one of the following:
\begin{itemize}
  \item \textbf{Measured:} a trend or pattern the data show directly.
  \item \textbf{Hypothesis:} a possible explanation the data are consistent with but don't prove.
  \item \textbf{Insufficient data:} a question the available records can't answer yet.
\end{itemize}

\subsection{Interactive Dashboard}
Results are presented in an interactive Streamlit dashboard built with pandas and Plotly. Code
and documentation are maintained on GitHub \citep{team23_repo}.
\draftnote{Dashboard lead: describe each page of the app and what a visitor can do there.}

% =====================================================================
\section{Results}\label{sec:results}
% =====================================================================
\subsection{Groundwater Levels}
\begin{figure}[t]
  \centering
  \figplaceholder[1.25in]{Figure 2 (to be added): time series of groundwater levels for wells with
  repeated measurements.}
  \caption{\draftnote{Caption to be written.}}
  \label{fig:levels}
\end{figure}
\draftnote{RQ1 (Fig.~\ref{fig:levels}).}

\subsection{Wells and Aquifer Units}
\draftnote{RQ2.}

\subsection{Cuba Case Study}
\draftnote{RQ3 and RQ4. Include the Rio Puerco at Cuba stream measurements as surface-water
context, not as evidence about groundwater.}

\subsection{Monitoring Coverage and Data Gaps}
\draftnote{RQ5.}

% =====================================================================
\section{Discussion}\label{sec:discussion}
% =====================================================================
\draftnote{What the results mean for Sandoval County and Cuba, and which indicators best
communicate groundwater stress (RQ6).}

% =====================================================================
\section{Limitations}\label{sec:limits}
% =====================================================================
The project proposal included three questions that this report does not answer, because the data
needed for them were not obtained during the project:
\begin{itemize}
  \item \textbf{Groundwater quality and salinity.} No groundwater-quality records were collected.
        The water-quality measurements in the USGS file all come from streams, mainly the Rio
        Grande below Cochiti Dam, far from Cuba.
  \item \textbf{Acequias.} No data on acequia locations or service areas were collected.
  \item \textbf{Groundwater and surface-water links.} Comparing the two requires repeated
        water-level measurements near a stream. The only stream records at Cuba are 10 visits to
        the Rio Puerco gage, and the nearest well in the water-level datasets is about 24 miles
        away.
\end{itemize}
\draftnote{Add the other data gaps and what the data can't support.}

% =====================================================================
\section{Conclusions}\label{sec:conclusions}
% =====================================================================
\draftnote{Summary of findings and recommended next steps.}

% =====================================================================
%  Back matter
% =====================================================================
\Needspace{9\baselineskip}
\section*{Data and Code Availability}
\begin{description}[style=nextline,leftmargin=1em,itemsep=1pt,font=\normalfont\bfseries]
  \item[Code and data] \href{https://github.com/pdjordj11/eng_220_aquifer_project}{\nolinkurl{github.com/pdjordj11/eng_220_aquifer_project}}
  \item[Interactive dashboard] \href{https://eng220aquiferproject.streamlit.app/}{\nolinkurl{eng220aquiferproject.streamlit.app}}
  \item[Documentation] \href{https://pdjordj11.github.io/eng_220_aquifer_project/}{\nolinkurl{pdjordj11.github.io/eng_220_aquifer_project}}
\end{description}

\section*{Author Contributions}
\begin{description}[leftmargin=0pt,labelsep=0.4em,itemsep=0pt,font=\normalfont\bfseries]
  \item[P.~Djordjevic:] documentation, final report, presentation, and video.
  \item[D.~Gonzales:] data acquisition: collected the datasets and wrote a research brief for
        each, plus a plotting guide comparing them.
  \item[J.~Wauneka:] data cleaning: wrote the MATLAB script that produces the cleaned
        groundwater-level and aquifer-test files.
  \item[M.~Watkins:] analysis code and visualizations.
  \item[M.~Yousef:] GitHub workflow and the Streamlit app.
\end{description}

\section*{Acknowledgments}
\draftnote{To be written.}

\bibliographystyle{plainnat}
\bibliography{references}

\end{document}
